\documentclass{article}
%%%%% PROGRAMMING PACKAGES
\usepackage{expl3,xparse}     %LaTeX3
\usepackage{mfirstuc}    %first letter uppercase with \makefirstuc (used in todo notes)
\usepackage{settobox}     %get box dimensions with \settoboxwidth (used in fit environment) => pretty stupid package.. could be avoided

%%%%% FIGURES, BOXES, COLORS
\usepackage{graphics,graphicx,trimclip,adjustbox}     %figures and box management
\usepackage{color,xcolor}     %to use colors (used in the todonotes)
    
%%%%% ENUMERATE
\usepackage{enumitem}   %to use "enumerate"
    \setlist[enumerate,1]{label = $\bullet$}   %default label at bullet points
    
%%%%% TODO NOTES
\usepackage[textwidth=16.3mm,textsize=scriptsize]{todonotes}   %to put notes everywhere
\newcommand*{\newcommentator}[2][red!20]{%
    \expandafter\newcommand\csname #2\endcsname[1]{\todo[color=#1]{{\bf\makefirstuc{#2}:} ##1}}%
    \expandafter\newcommand\csname big#2\endcsname[1]{\todo[inline,color=#1]{{\bf\makefirstuc{#2}:} ##1}}}
\newcommentator{pablo}

%%%%% MY HYPERLINK AND REF/LABEL MACROS
\usepackage{mypreamble/linkref}

%%%%% MATH PACKAGES
\usepackage{amssymb,amsmath}     %math original package
    \numberwithin{equation}{section}
\usepackage{mathtools}     %improves AMSmath
\usepackage{stackrel}     %better stacking of math symbols
\usepackage{mleftright}     %defines \mleft and \mright
    \mleftright     %redefines \left as \mleft and \right as \mright
\usepackage{pgfcore}
\usepackage{tikz-cd}     %advanced and rigorous diagrams
    
%%%%% MY SPECIAL PACKAGES
\usepackage{mypreamble/mathenv}     %improves on AMSmath's equation
\usepackage{mypreamble/thmenv}     %replaces the AMSthm package

%%%%% MATH MACROS AND FONT PACKAGES
%\usepackage{fontspec}
\usepackage{mypreamble/mathcro}

%%%%% HYPERREF PACKAGE
\makeatletter
\pretocmd\refstepcounter{%
 \zref@setcurrent{counter}{#1}}{}{}
\makeatother
\usepackage[pdfencoding=auto]{hyperref}   %hyperlinks parts of the pdf
    \hypersetup{
        bookmarksopen=true,   %expends tree of bookmarks in Acrobat pdf reader
        bookmarksnumbered=true,  %show numbers of sections in bookmarks
        linktoc=all,   %both section and page numbers are links
        hidelinks   %links (color+border) are hidden
    }
\providecommand*{\definitionautorefname}{Definition}
\providecommand*{\notationautorefname}{Notation}
\providecommand*{\conventionautorefname}{Convention}
\providecommand*{\conjectureautorefname}{Conjecture}
\providecommand*{\assumptionautorefname}{Assumption}
\providecommand*{\remarkautorefname}{Remark}
\providecommand*{\observationautorefname}{Observation}
\providecommand*{\lemmaautorefname}{Lemma}
\providecommand*{\propositionautorefname}{Proposition}
\providecommand*{\corollaryautorefname}{Corollary}
\providecommand*{\exampleautorefname}{Example}
\providecommand*{\exerciseautorefname}{Exercise}
\providecommand*{\openquestionautorefname}{Open Question}

%%%%% OUTPUT LAYOUT PACKAGES
\usepackage{geometry}     %to set margins
\geometry{
    a4paper,
    total={170mm,257mm},
    left=20mm,
    top=20mm,
    marginparwidth=20mm,
}
%% Clearpage before sections but not before TOC
\let\oldsection\section
\renewcommand\section{\clearpage\oldsection}
\makeatletter
\renewcommand\tableofcontents{%
    \oldsection*{\contentsname\@mkboth{\MakeUppercase\contentsname}{\MakeUppercase\contentsname}}\@starttoc{toc}%
}
\makeatother    %ensures the TOC stays in front page

%%%%% TYPESET DEBUGGING
%\usepackage{showframe}     %to make margins appear
%\usepackage{lua-visual-debug}     %to make all boxes appear
%\overfullrule=1pt   %creates 1pt wide black bars next to any overfull hbox

%%%%% QUIVER SETUP
\usetikzlibrary{calc,decorations.pathmorphing}
% A TikZ style for curved arrows of a fixed height, due to AndréC.
\tikzset{curve/.style={settings={#1},to path={(\tikztostart)
    .. controls ($(\tikztostart)!\pv{pos}!(\tikztotarget)!\pv{height}!270:(\tikztotarget)$)
    and ($(\tikztostart)!1-\pv{pos}!(\tikztotarget)!\pv{height}!270:(\tikztotarget)$)
    .. (\tikztotarget)\tikztonodes}},
    settings/.code={\tikzset{quiver/.cd,#1}
        \def\pv##1{\pgfkeysvalueof{/tikz/quiver/##1}}},
    quiver/.cd,pos/.initial=0.35,height/.initial=0}

% TikZ arrowhead/tail styles.
\tikzset{tail reversed/.code={\pgfsetarrowsstart{tikzcd to}}}
\tikzset{2tail/.code={\pgfsetarrowsstart{Implies[reversed]}}}
\tikzset{2tail reversed/.code={\pgfsetarrowsstart{Implies}}}
% TikZ arrow styles.
\tikzset{no body/.style={/tikz/dash pattern=on 0 off 1mm}}
\usepackage{mathdots}
\usepackage{bookmark,bbm}

\newcommand{\StRing}{\mathsf{StRing}}

%%%%% TITLE PAGE
\title{Gestalten}
\author{Pablo Bustillo Vazquez}

\begin{document}

\maketitle
\tableofcontents
\newpage

\subsection{Dualizability Composes}

Let $A \in \StRing$ and $B_n \inset \EAlg{\infty}(A_{n+1})$.
Then we get an adjunction
% https://q.uiver.app/#q=WzAsMixbMCwwLCJBX3tuKzF9Il0sWzIsMCwiXFxNb2Rfe0Jfbn0oQV97bisxfSkiXSxbMCwxLCJmXiogPSAtIFxcb3RpbWVzIEJfbiIsMCx7ImN1cnZlIjotMn1dLFsxLDAsImZfKiIsMCx7ImN1cnZlIjotMn1dLFsyLDMsIiIsMCx7ImxldmVsIjoxLCJzdHlsZSI6eyJuYW1lIjoiYWRqdW5jdGlvbiJ9fV1d
\[\begin{tikzcd}
	{A_{n+1}} && {\Mod_{B_n}(A_{n+1})}
	\arrow[""{name=0, anchor=center, inner sep=0}, "{f^* = - \otimes B_n}", curve={height=-12pt}, from=1-1, to=1-3]
	\arrow[""{name=1, anchor=center, inner sep=0}, "{f_*}", curve={height=-12pt}, from=1-3, to=1-1]
	\arrow["\dashv"{anchor=center, rotate=-90}, draw=none, from=0, to=1]
\end{tikzcd}\]

When does this adjunction preserves/reflects notions of dualizability? How about in more/less affine cases?

\subsection{Rigidity}

We fix $A \in \StRing$.
Note that for $A$ 1-affine and $B_1 \inset \EAlg{\infty}(A_{2})$, rigidity of the $B_1$ is equivalent to $f\in A \to B$ and $\Delta_f \in A \to B$ being $1$-prim (Cor. $4.45$ in \cite{Ram24a}).

\begin{definition}[(Local) Rigidity]
    Let $B_n \inset \EAlg{\infty}(A_{n+1})$ and the induced $f \in A \to B$.
    We say that $B_n$ is \emph{rigid} if $f$ is $n$-proper (I think the higher $m$ conditions are redundant here).
    We could define a weaker notion of \emph{$k$-local rigidity} by only asking that $\Delta^{(k)}_f, \Delta^{(k+1)}_f, \ldots$ are $n$-prim and that $B_n$ is dualizable as a $\int_{S^{-1}} B_n$-module, $\int_{S^{0}} B_n$-module, \ldots, $\int_{S^{k-1}} B_n$-module.
\end{definition}

Then hopefully, local-rigidity implies rigidity (which is $-1$-local rigidity) BUT one needs to find a non-affine definition to make this interesting. 

\section{Dualizability}

\begin{definition}[$n$-prim]
    $f\in A \to B$ is $n$-prim if, for all $m \geq n+1$, we have 
% https://q.uiver.app/#q=WzAsMixbMiwwLCIoQi9BKV9tIl0sWzAsMCwiXFxtYXRoYmJtezF9X20iXSxbMSwwLCJmX3ttLTF9Xlxcc3RhciIsMCx7ImN1cnZlIjotMn1dLFswLDEsImZfe20tMSwgXFxzdGFyfSIsMCx7ImN1cnZlIjotMiwic3R5bGUiOnsiYm9keSI6eyJuYW1lIjoiZGFzaGVkIn19fV0sWzIsMywiIiwyLHsibGV2ZWwiOjEsInN0eWxlIjp7Im5hbWUiOiJhZGp1bmN0aW9uIn19XV0=
\[\begin{tikzcd}
	{\mathbbm{1}_m} && {(B/A)_m}
	\arrow[""{name=0, anchor=center, inner sep=0}, "{f_{m-1}^\star}", curve={height=-12pt}, from=1-1, to=1-3]
	\arrow[""{name=1, anchor=center, inner sep=0}, "{f_{m-1, \star}}", curve={height=-12pt}, dashed, from=1-3, to=1-1]
	\arrow["\dashv"{anchor=center, rotate=-90}, draw=none, from=0, to=1]
\end{tikzcd}\]
\end{definition}


\begin{definition}[$n$-proper]
    $f\in A \to B$ is $n$-proper if $f$, $\Delta_f$, $\Delta_{\Delta_f}$, \ldots are all $n$-prim.
\end{definition}

\begin{definition}[$n$-suave]
    $f\in A \to B$ is $n$-suave if, for all $m \geq n+1$, we have 
% https://q.uiver.app/#q=WzAsMyxbMCwwLCIoQi9BKV9tXiEiXSxbNCwwLCIoQi9BKV9tIl0sWzIsMCwiXFxtYXRoYmJtezF9X20iXSxbMCwyLCJmX3ttLTEsICF9IiwwLHsiY3VydmUiOi0yfV0sWzIsMCwiZl97bS0xfV4hIiwwLHsiY3VydmUiOi0yLCJzdHlsZSI6eyJib2R5Ijp7Im5hbWUiOiJkYXNoZWQifX19XSxbMiwxLCJmX3ttLTF9Xlxcc3RhciIsMix7ImN1cnZlIjoyfV0sWzEsMiwiZl97bS0xLCBcXCN9IiwyLHsiY3VydmUiOjIsInN0eWxlIjp7ImJvZHkiOnsibmFtZSI6ImRhc2hlZCJ9fX1dLFszLDQsIiIsMix7ImxldmVsIjoxLCJzdHlsZSI6eyJuYW1lIjoiYWRqdW5jdGlvbiJ9fV0sWzYsNSwiIiwwLHsibGV2ZWwiOjEsInN0eWxlIjp7Im5hbWUiOiJhZGp1bmN0aW9uIn19XV0=
\[\begin{tikzcd}
	{(B/A)_m^!} && {\mathbbm{1}_m} && {(B/A)_m}
	\arrow[""{name=0, anchor=center, inner sep=0}, "{f_{m-1, !}}", curve={height=-12pt}, from=1-1, to=1-3]
	\arrow[""{name=1, anchor=center, inner sep=0}, "{f_{m-1}^!}", curve={height=-12pt}, dashed, from=1-3, to=1-1]
	\arrow[""{name=2, anchor=center, inner sep=0}, "{f_{m-1}^\star}"', curve={height=12pt}, from=1-3, to=1-5]
	\arrow[""{name=3, anchor=center, inner sep=0}, "{f_{m-1, \#}}"', curve={height=12pt}, dashed, from=1-5, to=1-3]
	\arrow["\dashv"{anchor=center, rotate=-90}, draw=none, from=0, to=1]
	\arrow["\dashv"{anchor=center, rotate=-90}, draw=none, from=3, to=2]
\end{tikzcd}\]
    in $A_{m+1}$, where $(B/A)_m^! \to \mathbbm{1}_m$ is the counit of the $f_{m, \#} \dashv f_{m}^\star$ adjunction in $A_{m+2}$.
\end{definition}


\begin{definition}[$n$-étale]
    $f\in A \to B$ is $n$-étale if $f$, $\Delta_f$, $\Delta_{\Delta_f}$, \ldots are all $n$-suave.
\end{definition}

\begin{remark}[Redundancy]
    For $m \geq n+1$, we have $(B/A)_m^! \simeq (B/A)_m$ and $f_{m}^\star \dashv f_{m, \#}$.
    In this range, we can denote $f_{m, \star} = f_{m, \#}$ as it lifts the external right adjoint (which always exists but isn't necessarily $A$-linear).
\end{remark}

\subsection{$n$-proper is useful to understand diagonals and tensoring}

\begin{remark}[Higher diagonals]
    We get functors $(B \otimes_A B)_{m+1} \to \End_{A_{m+2}} ((B/A)_{m+1})$.
\end{remark}

\begin{lemma}
    If $f\in A \to B$ is $n$-prim, then $(B/A)_m \otimes_{A_{m+1}} (B/A)_m \simeq (B \otimes_A B /A)_m$ for $m \geq n$.
\end{lemma}
\begin{proof}
    We can internally describe $(B/A)_m = f_{m, \star} \circ f_{m}^\star$.
    A Eckmann-Hilton argument then shows that $(B/A)_m \otimes_{A_{m+1}} (B/A)_m \simeq \End({(B/A)_{m+1} \otimes_{A_{m}} (B/A)_{m+1}})$.
    Hence ``spectrification'' only changes $(B/A)_m \otimes_{A_{m+1}} (B/A)_m$ for $m < n$.
\end{proof}
\begin{corollary}
    If $f \in A \to B$ and $\Delta_f$ is $n$-prim, then $(B/A)_m$ is its own dual for $m \geq n+1$.
\end{corollary}
\begin{proof}
    Similar to showin that rigid implies frobenius implies self-dualizable.
    The Beck-Chevalley condition comes from $\Delta_f$ being $n$-prim (the adjoint is $B \otimes_A B$-linear).
\end{proof}
\pablo{the result above is nice, but etaleness will give an easier way to describe duality I think}

\begin{conjecture}
    If $f\in A \to B$ is $n$-proper, then the map $\int_{S^k} (B/A)_m \simeq S^k \otimes (B/A)_m \to (B/A)_m$ induced by the terminal $S^k \to \{\bullet\}$ is $(\Delta^{(k+1)}/A)_m$.
\end{conjecture}
\begin{proof}
    If $f$ is $n$-proper, then $A \to B \otimes_A B$ is $n$-prim.
    I \textbf{think} this should imply that tensoring internally in $(B \otimes_A B)_{m+1}$ is the same as tensoring internally in $A_{m+1}$ over $(B \otimes_A B/A)_{m+1}$.
    Then this should imply the result by induction on $k$.
\end{proof}

\bibliographystyle{alpha}
\bibliography{../mypreamble/jabref}

\end{document}